% ============================================================
%  TERESA Z1 — Diario di progresso e TODO list
%  Aggiornato: 2026-03-15
% ============================================================
\documentclass[a4paper,11pt]{article}

\usepackage[T1]{fontenc}
\usepackage[utf8]{inputenc}
\usepackage[italian]{babel}
\usepackage{geometry}
\usepackage{enumitem}
\usepackage{hyperref}
\usepackage{titlesec}
\usepackage{xcolor}
\usepackage{booktabs}
\usepackage{longtable}
\usepackage{fancyhdr}
\usepackage{mdframed}
\usepackage{listings}
\usepackage{parskip}
\usepackage{amssymb}

\geometry{margin=2.5cm}

% ── Colori ─────────────────────────────────────────────────
\definecolor{donegreen}{RGB}{0,140,0}
\definecolor{todoorange}{RGB}{200,80,0}
\definecolor{codeblue}{RGB}{30,60,150}
\definecolor{codegray}{RGB}{245,245,245}
\definecolor{warnyellow}{RGB}{180,130,0}

% ── Listings style ─────────────────────────────────────────
\lstset{
  basicstyle=\ttfamily\small,
  backgroundcolor=\color{codegray},
  breaklines=true,
  frame=single,
  rulecolor=\color{gray!40},
  numbers=none,
}

% ── Header / footer ───────────────────────────────────────
\pagestyle{fancy}
\fancyhf{}
\lhead{\textbf{TERESA Z1} — Braccio robotico Unitree Z1}
\rhead{Aggiornato: 2026-03-15}
\cfoot{\thepage}

% ── Ambienti custom ───────────────────────────────────────
\newmdenv[linecolor=donegreen!60, backgroundcolor=donegreen!5,
          innerleftmargin=8pt, innerrightmargin=8pt]{donebox}
\newmdenv[linecolor=todoorange!60, backgroundcolor=todoorange!5,
          innerleftmargin=8pt, innerrightmargin=8pt]{todobox}
\newmdenv[linecolor=warnyellow!80, backgroundcolor=warnyellow!5,
          innerleftmargin=8pt, innerrightmargin=8pt]{warnbox}

% ─────────────────────────────────────────────────────────
\begin{document}

\begin{center}
  {\LARGE \textbf{TERESA — Braccio Unitree Z1}}\\[4pt]
  {\large Diario di sviluppo, TODO list e guida di avvio}\\[2pt]
  {\small Branch: \texttt{FSM\_main} \quad|\quad Repo:
    \texttt{TERESA\_ws/src/z1\_vision}}\\[2pt]
  {\small Ultimo aggiornamento: \textbf{15 Marzo 2026}}
\end{center}
\vspace{0.5em}
\hrule
\vspace{1em}

% ============================================================
\section{Panoramica del sistema}
% ============================================================

Il sistema controlla il braccio \textbf{Unitree Z1} montato sul robot TERESA per
avvicinarsi in modo autonomo al torso di un paziente.
Lo stack completo è composto da:

\begin{itemize}
  \item \textbf{Percezione}: YOLO11 pose estimation + Kalman 3D + lock tracker
  \item \textbf{Pianificazione}: Workspace checker (Pinocchio + L-BFGS-B)
  \item \textbf{Controllo posizione}: IK Jacobiana pseudo-inversa smorzata + JointTrajectoryController
  \item \textbf{Controllo forza}: Impedance controller in spazio cartesiano (500 Hz, torque control)
  \item \textbf{Sicurezza}: keyboard safety node, stati HOMING/EMERGENCY, workspace clipping
  \item \textbf{Orchestrazione}: FSM a 11 stati, launch files separati per controllo e percezione
\end{itemize}

% ============================================================
\section{Architettura FSM (stato attuale)}
% ============================================================

\subsection*{Stati della FSM (\texttt{z1\_FSM.py})}

\begin{center}
\begin{tabular}{lp{10cm}}
\toprule
\textbf{Stato} & \textbf{Descrizione} \\
\midrule
\texttt{HOMING}              & Stato iniziale al boot. Invia home pose via IK/JTC, poi va in WAITING \\
\texttt{WAITING}             & Aspetta target LOCKED fresco ($<$0.5s) dal tracker YOLO \\
\texttt{CHECKING\_WORKSPACE} & Verifica raggiungibilità target (thread non bloccante, Pinocchio) \\
\texttt{APPROACHING}         & Invia goal IK verso il target (latchato al momento del check) \\
\texttt{WAIT\_IK\_DONE}      & Attende \texttt{/ik\_done=True} dal nodo IK \\
\texttt{SWITCHING\_TO\_TORQUE} & Switch JTC $\to$ torque\_controller via safe\_controller\_switch \\
\texttt{IMPEDANCE\_RUNNING}  & Impedance controller attivo: approach $\to$ hold $\to$ retract \\
\texttt{SWITCHING\_TO\_JTC}  & Switch torque\_controller $\to$ JTC \\
\texttt{EMERGENCY\_SWITCHING}& Forza switch a JTC se torque attivo, poi va in HOMING \\
\texttt{EMERGENCY}           & Stato bloccante, aspetta solo comando \texttt{reset} da tastiera \\
\texttt{FAULT}               & Errore fatale (switch fallito), nodo bloccato \\
\bottomrule
\end{tabular}
\end{center}

\subsection*{Topic principali}
\begin{center}
\begin{tabular}{lll}
\toprule
\textbf{Topic} & \textbf{Tipo} & \textbf{Descrizione} \\
\midrule
\texttt{/torso\_target\_ee\_locked}     & PoseStamped & Target EE quando tracker LOCKED \\
\texttt{/torso\_target\_ee}             & PoseStamped & Target EE interpolato continuo \\
\texttt{/ik\_goal\_pose}                & PoseStamped & Goal IK da FSM \\
\texttt{/ik\_enable}                    & Bool        & Abilita nodo IK \\
\texttt{/ik\_done}                      & Bool        & Conferma fine movimento JTC \\
\texttt{/impedance\_enable}             & Bool        & Abilita impedance controller \\
\texttt{/impedance\_done}              & Bool        & Fine ciclo impedance \\
\texttt{/target\_out\_of\_workspace}    & Bool        & Target clippato dal workspace checker \\
\texttt{/z1\_fsm/state}                 & String      & Stato corrente FSM \\
\texttt{/z1\_keyboard\_cmd}             & String      & Comandi da tastiera (home/emergency/reset) \\
\texttt{/torso\_surface\_frame}         & PoseStamped & Piano superficie torso (PCA) \\
\texttt{/surface\_signed\_distance}     & Float32     & Distanza TCP--superficie \\
\bottomrule
\end{tabular}
\end{center}

% ============================================================
\section{DONE — Funzionalità implementate}
% ============================================================

\begin{donebox}
\textbf{\color{donegreen} Funzionalità completate al 2026-03-15}

\subsection*{Pipeline base (commit precedenti)}
\begin{itemize}[leftmargin=1.5em]
  \item[\checkmark] Nodo \texttt{z1\_yolo\_torso\_tracker}: YOLO11 pose estimation,
        Kalman 3D [x,y,z,vx,vy,vz] con dt adattivo, FSM IDLE/ESTIMATING/LOCKED,
        lock su varianza 20 frame + 5 check consecutivi
  \item[\checkmark] Nodo \texttt{z1\_ik\_to\_jtc}: IK Jacobiana pseudo-inversa smorzata
        (frame LOCAL), smoothstep quintico ($10t^3-15t^4+6t^5$), joint unwrap anti-flip,
        interfaccia action JTC
  \item[\checkmark] IK funzionante sul robot reale (commit \texttt{0c5945d})
  \item[\checkmark] JTC funzionante e collegato alla FSM (commit \texttt{2eb6ae6})
  \item[\checkmark] Nodo \texttt{realsense\_surface\_node}: stima piano torso da depth
        ROI con PCA, pubblica frame superficie e distanza segnata TCP--piano
\end{itemize}

\subsection*{Impedance controller (commit \texttt{553ff7d}, 2026-03-15)}
\begin{itemize}[leftmargin=1.5em]
  \item[\checkmark] \texttt{impedance\_controller\_realsense.py}: mini state machine
        interna APPROACH (20 cm lungo normale) $\to$ HOLD (10s) $\to$ RETRACT (20 cm ritorno)
  \item[\checkmark] Interfaccia \texttt{/impedance\_enable} (Bool) / \texttt{/impedance\_done} (Bool)
  \item[\checkmark] Safe startup 3s, compensazione dinamica (gravity+coriolis),
        scale adattivo J2 in funzione del reach XY, Jacobiano LOCAL\_WORLD\_ALIGNED
  \item[\checkmark] \texttt{safe\_controller\_switch.py}: convertito in nodo persistente con
        servizi ROS2 \\\texttt{/safe\_switch/to\_torque} e \texttt{/safe\_switch/to\_jtc}
  \item[\checkmark] FSM estesa a 7 stati con ciclo JTC $\to$ torque $\to$ JTC
  \item[\checkmark] Vecchi file spostati in \texttt{Old/}:
        \texttt{trajectory\_manager.py}, \texttt{z1\_FSM\_old.py},
       \\ \texttt{z1\_ik\_to\_jtc\_old.py}, \texttt{z1\_ik\_to\_jtc\_standalone.py}
\end{itemize}

\subsection*{Workspace checker (commit \texttt{dd6c4b7}, 2026-03-15)}
\begin{itemize}[leftmargin=1.5em]
  \item[\checkmark] \texttt{workspace\_checker.py}: calcolo raggiungibilità massima
        con Pinocchio + scipy L-BFGS-B, clip target a (max\_reach - 30 cm)
  \item[\checkmark] Stato \texttt{CHECKING\_WORKSPACE} nella FSM (thread non bloccante,
        ThreadPoolExecutor)
  \item[\checkmark] Target latchato al lancio del thread: il robot completa sempre
        il movimento anche se il tracker perde il lock
  \item[\checkmark] Topic \texttt{/target\_out\_of\_workspace} (Bool) per integrazione futura
        con Spot
  \item[\checkmark] Rimossi gli abort per target non fresco in APPROACHING e WAIT\_IK\_DONE
\end{itemize}

\subsection*{HOMING + EMERGENCY + keyboard safety (commit \texttt{22c15a2}, 2026-03-15)}
\begin{itemize}[leftmargin=1.5em]
  \item[\checkmark] Stato iniziale FSM: \texttt{HOMING} (porta il braccio in
        \texttt{home\_position} all'avvio, poi WAITING)
  \item[\checkmark] Stato \texttt{EMERGENCY\_SWITCHING}: forza switch JTC se torque attivo,
        poi torna in HOMING
  \item[\checkmark] Stato \texttt{EMERGENCY}: bloccante, aspetta solo \texttt{reset} da tastiera
  \item[\checkmark] Nuovo nodo \texttt{z1\_keyboard\_safety.py}: raw terminal mode (termios),
        tasti H=Home, ESC=Emergency, R=Reset, Q=Quit; refresh stato 2Hz
  \item[\checkmark] Parametri YAML per home pose: \texttt{home\_position},
        \texttt{home\_orientation}
  \item[\checkmark] Launch file \texttt{z1\_control.launch.py}: safe\_controller\_switch +
        z1\_ik\_to\_jtc + impedance (t=0s), FSM (t=5s delay)
  \item[\checkmark] Launch file \texttt{z1\_perception.launch.py}: YOLO tracker +
        surface node
\end{itemize}
\end{donebox}

% ============================================================
\section{TODO — Da completare}
% ============================================================

\begin{todobox}
\textbf{\color{todoorange} Attività in sospeso}

\subsection*{Alta priorità}
\begin{itemize}[leftmargin=1.5em]
  \item[$\square$] \textbf{Calibrare la home position reale}: il valore attuale in
        \texttt{z1\_fsm\_params.yaml} è un placeholder
        (\texttt{[0.3, 0.0, 0.3]}, quaternione identità). Verificare con il robot
        reale qual è la posa di riposo sicura e aggiornare il YAML.
  \item[$\square$] \textbf{Testare la sequenza completa sul robot reale}:
        HOMING $\to$ WAITING $\to$ CHECKING\_WORKSPACE $\to$ APPROACHING $\to$
        WAIT\_IK\_DONE $\to$ \\ SWITCHING\_TO\_TORQUE $\to$ IMPEDANCE\_RUNNING
        $\to$ SWITCHING\_TO\_JTC $\to$ WAITING
  \item[$\square$] \textbf{Validare il workspace checker} con target reali:
        verificare che il clip a (max\_reach - 30 cm) sia corretto e che la
        posa clippata sia effettivamente raggiungibile
  \item[$\square$] \textbf{Testare z1\_keyboard\_safety} sul robot reale:
        verificare che H, ESC e R funzionino in tutti gli stati FSM (inclusi
        gli stati torque)
\end{itemize}

\subsection*{Media priorità}
\begin{itemize}[leftmargin=1.5em]
  \item[$\square$] \textbf{Integrazione con Spot / base mobile}: usare
        \texttt{/target\_out\_of\_workspace} per comandare alla base di
        avvicinarsi quando il target è fuori workspace
  \item[$\square$] \textbf{Timeout e recovery in IMPEDANCE\_RUNNING}: aggiungere
        timeout di sicurezza per il caso in cui il ciclo impedance non
        termini (es. \texttt{/impedance\_done} non arriva mai)
  \item[$\square$] \textbf{Tuning parametri impedance controller}: guadagni
        cartesiani, distanza di approach, durata hold (10s configurabile?)
  \item[$\square$] \textbf{z1\_keyboard\_safety nel launch file}: valutare se
        aggiungere un modo per avviarlo automaticamente in un terminale
        dedicato (es. \texttt{xterm} o \texttt{gnome-terminal} nel launch file)
\end{itemize}

\subsection*{Bassa priorità / futuri miglioramenti}
\begin{itemize}[leftmargin=1.5em]
  \item[$\square$] \textbf{Logging e monitoraggio}: aggiungere ROS bag recording
        automatico al launch con i topic principali
  \item[$\square$] \textbf{Visualizzazione RViz}: marker per il target, workspace
        reachability, piano superficie
  \item[$\square$] \textbf{Multi-persona}: gestione del caso in cui YOLO rileva
        più persone (attualmente prende la prima/più vicina)
  \item[$\square$] \textbf{Calibrazione estrinseci camera-braccio}: verificare e
        documentare la trasformazione TF tra frame camera e frame base Z1
\end{itemize}
\end{todobox}

% ============================================================
\section{Sequenza di avvio}
% ============================================================

\begin{warnbox}
\textbf{Prerequisito}: il workspace ROS2 deve essere compilato e sourced.\\
\texttt{cd \textasciitilde/Documents/GIT\_Repositories/TERESA\_ws}\\
\texttt{colcon build --packages-select z1\_vision \&\& source install/setup.bash}
\end{warnbox}

\subsection*{Terminale 1 — Hardware robot + camera + TF}

\begin{lstlisting}
# Avvia driver hardware Unitree Z1, RealSense D435, TF statiche
ros2 launch z1_vision z1_realsense.launch.py
\end{lstlisting}

\textit{Attende che il JointTrajectoryController sia attivo e la camera streaming.}

\subsection*{Terminale 2 — Percezione (YOLO + surface)}

\begin{lstlisting}
ros2 launch z1_vision z1_perception.launch.py

# Opzionale: solo tracker senza surface node
# ros2 launch z1_vision z1_perception.launch.py use_surface_node:=false
\end{lstlisting}

\subsection*{Terminale 3 — Stack di controllo}

\begin{lstlisting}
# Avvia: safe_controller_switch + z1_ik_to_jtc + impedance (t=0s)
#        z1_FSM dopo 5s (primo stato: HOMING)
ros2 launch z1_vision z1_control.launch.py

# Opzionale: aumentare il delay se serve piu' tempo
# ros2 launch z1_vision z1_control.launch.py fsm_delay:=8.0
\end{lstlisting}

\textit{Attende il completamento dell'HOMING prima di passare in WAITING.}

\subsection*{Terminale 4 — Keyboard safety (interattivo, separato)}

\begin{lstlisting}
ros2 run z1_vision z1_keyboard_safety
\end{lstlisting}

\begin{center}
\begin{tabular}{ll}
\toprule
\textbf{Tasto} & \textbf{Azione} \\
\midrule
\texttt{H}   & Invia braccio alla home position (qualunque stato) \\
\texttt{ESC} & Emergenza: stop immediato, poi HOMING \\
\texttt{R}   & Reset dallo stato EMERGENCY \\
\texttt{Q}   & Chiude il nodo keyboard safety \\
\bottomrule
\end{tabular}
\end{center}

\subsection*{Monitoraggio (terminali aggiuntivi opzionali)}

\begin{lstlisting}
# Stato FSM in tempo reale
ros2 topic echo /z1_fsm/state

# Target YOLO
ros2 topic echo /torso_target_ee_locked

# Distanza TCP-superficie
ros2 topic echo /surface_signed_distance

# Target fuori workspace
ros2 topic echo /target_out_of_workspace

# FK end-effector
ros2 run z1_vision fk_reader
\end{lstlisting}

% ============================================================
\section{File principali}
% ============================================================

\begin{center}
\begin{longtable}{lp{9cm}}
\toprule
\textbf{File} & \textbf{Ruolo} \\
\midrule
\endhead
\texttt{z1\_FSM.py}                        & FSM principale (11 stati) \\
\texttt{z1\_ik\_to\_jtc.py}                & IK Pinocchio $\to$ JTC action \\
\texttt{impedance\_controller\_realsense.py}& Controllo impedenza cartesiana 500 Hz \\
\texttt{safe\_controller\_switch.py}        & Switch JTC $\leftrightarrow$ torque (servizi ROS2) \\
\texttt{workspace\_checker.py}             & Raggiungibilità workspace (Pinocchio + scipy) \\
\texttt{z1\_yolo\_torso\_tracker.py}        & YOLO11 + Kalman 3D + lock \\
\texttt{realsense\_surface\_node.py}        & Stima piano torso (PCA depth ROI) \\
\texttt{z1\_keyboard\_safety.py}           & Safety node interattivo da terminale \\
\texttt{kalman\_filter.py}                 & Filtro Kalman 3D [x,y,z,vx,vy,vz] \\
\texttt{fk\_reader.py}                     & Utility: FK end-effector da joint\_states \\
\midrule
\texttt{config/z1\_fsm\_params.yaml}        & Parametri FSM (home pose, topic, workspace) \\
\texttt{config/z1\_ik\_jtc\_params.yaml}    & Parametri IK/JTC (URDF, tolleranze, traiettoria) \\
\texttt{config/z1\_yolo\_torso\_params.yaml}& Parametri YOLO/Kalman/lock \\
\texttt{config/impedance\_control\_params.yaml} & Parametri impedance (guadagni, timing) \\
\texttt{config/surface\_params.yaml}        & Parametri surface node \\
\midrule
\texttt{launch/z1\_control.launch.py}       & Avvia stack controllo (switch + IK + impedance + FSM) \\
\texttt{launch/z1\_perception.launch.py}    & Avvia stack percezione (YOLO + surface) \\
\texttt{launch/z1\_realsense.launch.py}     & Avvia HW robot + camera + TF \\
\bottomrule
\end{longtable}
\end{center}

% ============================================================
\section{Note tecniche}
% ============================================================

\begin{itemize}
  \item \textbf{IK}: Jacobiana pseudo-inversa smorzata (\emph{damped least squares}),
        frame LOCAL, interpolazione smoothstep quintica (zero velocità a inizio/fine),
        joint unwrap con \texttt{\_make\_target\_near} per evitare rotazioni $>|\pi|$
  \item \textbf{Kalman}: dt adattivo dal timestamp immagine,
        \texttt{vel\_damping = 0.9}, covarianza noise configurabile da YAML
  \item \textbf{Lock tracker}: varianza soglia su 20 frame + 5 check consecutivi sotto soglia
  \item \textbf{Impedance}: 500 Hz, PID cartesiano posizione + joint-space polso (J4--J6),
        compensazione gravity+coriolis, Jacobiano LOCAL\_WORLD\_ALIGNED,
        scale adattivo J2 in funzione del reach XY
  \item \textbf{URDF path} (hardcoded): \\
        \texttt{/home/andrea/Ros2\_repositories/unitree\_z1\_ws/install/}\\
        \texttt{z1\_description/share/z1\_description/urdf/z1.urdf}
  \item \textbf{Branch corrente}: \texttt{FSM\_main}
        (fare merge su \texttt{main} dopo i test sul robot reale)
\end{itemize}

\end{document}
